%----------------------------------------------------------------------------------------
%	CHAPTER 3
%----------------------------------------------------------------------------------------

\chapter{Projekt oraz implementacja interfejsu SSVEP}

W niniejszym rozdziale przedstawiono główne założenia projektowe oraz ogólną koncepcję implementacji systemu interfejsu mózg-komputer opartego na zjawisku wzrokowych potencjałów wywołanych stanu ustalonego. Celem projektu jest opracowanie rozwiązania umożliwiającego sterowanie kursorem myszy za pomocą odpowiedzi neuronalnych generowanych w odpowiedzi na migoczące bodźce wzrokowe. Konstrukcja systemu musi zapewniać nie tylko funkcjonalność pozwalającą na wykonywanie podstawowych komend ruchu kursora, ale również powtarzalność wyników. Z tego względu nacisk położono na równowagę pomiędzy prostotą konfiguracji a jakością akwizycji sygnału EEG.

\begin{figure}[!ht]
    \centering
    \includegraphics[width=0.5\textwidth]{diagram_moj123123.png}
    \caption{Uproszczony schemat stanowiska, który zostanie użyty w eksperymencie.}
    \label{fig:schemat_moj}
\end{figure}

\subsection*{Założenia funkcjonalne systemu} 

\begin{itemize}
    \item System umożliwia sterowanie kursorem w czterech kierunkach oraz generowanie kliknięcia
    \item Użytkownik wybiera komendę poprzez skupienie wzroku na jednym z pięciu migających bodźców
    \item System dokonuje analizy sygnału w oknach czasowych o stałej długości
    \item Interfejs pozwala na konfigurację parametrów migających pól
\end{itemize}

\subsection*{Założenia pozafunkcjonalne systemu} 

\begin{itemize}
    \item Stabilność działania systemu nie zależy od jakości łącza Bluetooth
    \item Moduł przetwarzania sygnału musi być odporny na zakłócenia w warunkach laboratoryjnych
    \item Aplikacja działa w sposób deterministyczny
    \item System działa w czasie rzeczywistym
\end{itemize}



Do rejestracji sygnału EEG zostanie zastosowany moduł OpenBCI Cyton \cite{openbci} oparty na 8-kanałowym wzmacniaczu Texas Instruments ADS1299, 24-bitowym przetworniku analogowo-cyfrowy oraz 8-bitowym mikroprocesorze PIC32MX250F128B firmy Microchip Technology. Urządzenie zostanie połączone bezprzewodowo z komputerem za pomocą Bluetooth 4.0. 

\begin{figure}[!ht]
    \centering
    \includegraphics[width=0.5\textwidth]{Cyton.jpg}
    \caption{Urządzenie, które zostanie użyte do przeprowadzenia eksperymentu. Źródło: \cite{openbci}}
    \label{fig:cyton}
\end{figure}

W roli modułu stymulującego zastosowany zostanie ekran LCD, na którym umieszczone będą pola migające z różnymi częstotliwościami. Autor zdecydował się na to rozwiązanie ze względu na łatwość implementacji wielu niezależnych bodźców. Nakładka programowa zostanie zaimplementowana jako bezramkowe, przezroczyste okno, działające zawsze na wierzchu i ignorujące wejście użytkownika. Zostanie ona napisana w języku Python przy użyciu biblioteki PyQt5. Zdecydowano się na wariant pięciu niezależnych migających pól, z których cztery odpowiadać będą za sterowanie kierunkowe, natomiast piąty będzie pełnił funkcję symulacji kliknięcia. Na rysunku 3.1 został przedstawiony wygląd nakładki.

\begin{figure}[!ht]
    \centering
    \includegraphics[width=0.8\textwidth]{ssvep_stym.PNG}
    \caption{Zrzut ekranu przedstawiający stymulator. Kolor jednego z pól został zmieniony na czerwony.}
    \label{fig:ssvep_stym}
\end{figure}

Nakładka programowa umożliwi edycję parametrów wizualnych, takich jak położenie, kolor, rozmiar oraz przezroczystość poszczególnych pól, co pozwoli dostosować ich wygląd do wymogów eksperymentu oraz indywidualnych preferencji badanego. Renderowanie wszystkich elementów odbędzie się w rytmie zależnym od częstotliwości odświeżania monitora, która będzie odczytywana dynamicznie za pomocą polecenia systemowego wmic. Aplikacja wyznacza z tej wartości wszystkie dozwolone częstotliwości migotania, ograniczając je wyłącznie do dzielników częstotliwości odświeżania monitora. Mechanizm ten wymusza synchronizację migotania z częstotliwością odświeżania matrycy LCD. Każdy bodziec przedstawiony będzie jako prostokąt, którego widoczność będzie przełączana na podstawie licznika klatek. Logika migotania zostanie oparta na porównaniu aktualnego numeru klatki z liczbą klatek przypadającą na półokres danej częstotliwości. Dzięki temu częstotliwość prezentacji bodźca będzie deterministyczna, a jej wartość będzie wynikać bezpośrednio z relacji między liczbą klatek na sekundę a wybraną częstotliwością migotania. W docelowej konfiguracji stymulatory będą prezentowane na osobnym monitorze o częstotliwości odświeżania równej 60~Hz. To odseparuje warstwę interakcji użytkownika z komputerem od warstwy pobudzenia wzrokowego.

\begin{figure}[!ht]
    \centering
    \includegraphics[width=0.8\textwidth]{opcje_ssvep.PNG}
    \caption{Okienko pozwalające na zmianę parametrów poszczególnych pól stymulatora.}
    \label{fig:ssvep_options}
\end{figure}


Drugi podsystem będzie realizować rejestrację oraz wstępne przetwarzanie sygnału EEG. Dane będą pobierane z trzech elektrod umieszczonych w punktach O1, O2, Oz, elektrody odniesienia umieszczonej na wyrostku sutkowatym oraz elektrody uziemiającej umieszczonej na wyrostku sutkowatym drugiego ucha. Program zostanie napisany w języku Python z użyciem bibliotek mne, BrainFlow, scipy oraz numpy. Wstępna obróbka obejmuje filtrację pasmową 0.5-30~Hz. Dzięki temu wyeliminuje się także częstotliwość związaną z instalacją elektryczną (pasmo 50~Hz). Na każdym kanale zostanie wykonana szybka transformata Fouriera, która przekształci sygnał z dziedziny czasu do dziedziny częstotliwości. Na podstawie otrzymanego widma wyznaczony zostanie estymator gęstości widmowej mocy. Z tak obliczonego PSD zostaną wyodrębniane pasma odpowiadające pięciu częstotliwościom stymulacji, po czym wartości mocy w tych zakresach będą sumowane. W celu zwiększenia czułości detekcji uwzględniane będą również pierwsze harmoniczne każdej częstotliwości bodźca.

Klasyfikacja będzie polegać na wyznaczeniu maksymalnej wartości PSD w paśmie $\pm 0.5$~Hz wokół każdej częstotliwości bodźca. Bodziec, którego częstotliwość odpowiada składowej z najwyższą mocą, zostanie uznany za aktualnie obserwowaną w odpowiedzi mózgowej użytkownika. Wynik klasyfikacji wykorzystywany będzie następnie do sterowania ruchem kursora lub generowania kliknięcia myszy w interfejsie użytkownika.




\begin{lstlisting}[language=Python, caption=Fragment kodu obliczającego FFT przy użyciu biblioteki numpy, label=alg:fft]

import numpy as np

def compute_fft(x, fs):
    import numpy as np
    N = len(x)
    X = np.fft.rfft(x)
    freqs = np.fft.rfftfreq(N, 1/fs)
    return freqs, np.abs(X)

\end{lstlisting}


\begin{lstlisting}[language=Python, caption=Fragment algorytmu PSD, label=alg:psd]
import numpy as np

def compute_psd(x, fs):
    N = len(x)
    X = np.fft.rfft(x)
    psd = (1/N) * (np.abs(X) ** 2)
    freqs = np.fft.rfftfreq(N, 1/fs)
    return freqs, psd

\end{lstlisting}


\begin{lstlisting}[language=Python, caption=Fragment algorytmu klasyfikatora, label=alg:klasyk]
import numpy as np

def compute_snr(freqs, psd, fb, bw=1.0):
    mask_sig = (freqs >= fb - 0.5) & (freqs <= fb + 0.5)
    mask_noise = (freqs >= fb - bw) & (freqs <= fb + bw) & (~mask_sig)
    sig = np.mean(psd[mask_sig])
    noise = np.mean(psd[mask_noise])
    return sig / noise

def classify_snr(x, fs, freqs_targets):
    freqs, psd = compute_psd(x, fs)
    scores = {f: compute_snr(freqs, psd, f) for f in freqs_targets}
    return max(scores, key=scores.get), scores
\end{lstlisting}



